% SEGMENT OLP-0662-S01
% Part: proof-theory
% Chapter: natural-deduction
% Section: grafting

\documentclass[../../../include/open-logic-section]{subfiles}

\begin{document}

\olfileid{pt}{nat}{gra}

\olsection{நிறுவல்களை ஒட்டுதல் \usetoken{P}{proof}}

\Log{N1c} மற்றும்~\Log{N1i} இல் உள்ள !!^{proof}s,
எடுகோள்கள்~$!B$ இலிருந்து !!{formula}s~$!A$ க்கான !!{proof}s ஆகும்.
$!B$ க்கும் தனியே !!a{proof} உள்ளது என்க. அப்போது $!B$
இலிருந்து $!A$ வுக்கான !!{proof} உடன் $!B$ இன் !!{proof} ஐ
எவ்வாறு இணைத்தால், $!B$ என்ற எடுகோளைப் பயன்படுத்தாத $!A$ இன்
!!a{proof} கிடைக்கும்?

ஒரு வழி, $!A$ இன் !!{proof} இல் உள்ள எடுகோள்~$!B$ ஐ
\Intro{\lif} பயன்படுத்தி விடுவிப்பது. அதன் விளைவாகக் கிடைக்கும்
$!B \lif !A$ இன் !!{proof} உடன் $!B$ இன் !!{proof} ஐ
இணைத்தால், பின்வருவது கிடைக்கும்:
\[
\AxiomC{$\Discharge{!B}{x}$}
\DeduceC{$!A$}
\DischargeRule{\Intro{\lif}}{x}
\UnaryInfC{$!B \lif !A$}
\AxiomC{}
\DeduceC{$!B$}
\RightLabel{\Elim{\lif}}
\BinaryInfC{$!A$}
\DisplayProof
\]
இது நேரடியான வழிதான்; ஆனால் தேவையற்ற சுற்றுவழியும் ஆகும்.
உடனே நீக்குவதற்காக மட்டும் $\lif$ ஐ ஏன் அறிமுகப்படுத்த வேண்டும்?
$!B \lif !A$ வழியாகச் செல்லும் இத்தகைய ``சுற்றுவழியை''
தவிர்க்க முடியும். (அவற்றை எப்போதும் தவிர்க்க முடியும் என்பதே
இயல்பாக்கத் தேற்றத்தின் உள்ளடக்கம்.)

% SEGMENT OLP-0662-S02
\Log{N1c} மற்றும்~\Log{N1i} இல், இந்த நேர்வில் அதை எளிதாகத்
தவிர்க்கலாம்: எடுகோள்கள்~$!B^x$ இருக்கும் இடங்களில்
$!B$ இன் முழு !!{proof} ஐப் பதிலிடலாம். ஒரு நிறுவலின் எடுகோள்களில்
மற்றொரு !!{proof} ஐ இவ்வாறு ``ஒட்டுதல்'' பயனுள்ளதாக இருப்பதால்,
அதை ஒரு வரையறையாகப் பதிவு செய்கிறோம்.

\begin{defn}
$\delta_1$ என்பது திறந்த எடுகோள்~$!B^x$ இலிருந்து $!A$ வுக்கான
\Log{N1c}-!!{proof} (அல்லது \Log{N1i}-!!{proof}) ஆகவும்,
$\delta_2$ என்பது $!B$ இன் \Log{N1c}-!!{proof}
(அல்லது \Log{N1i}-!!{proof}) ஆகவும் இருந்தால்,
$\Subst{\delta_1}{\delta_2}{!B^x}$ என்பது $\delta_1$ இல்
விடுவிக்கப்படாமல் உள்ள $!B^x$ இன் ஒவ்வொரு நிகழ்வையும்
$\delta_2$ ஆல் பதிலிட்டதன் விளைவு.
\end{defn}

$\delta_1$ மற்றும் $\delta_2$ இல் வேறுபட்ட எடுகோள்
!!{formula}s ஒரே குறியீட்டால் குறிக்கப்படாமலும், $\delta_2$
இன் எந்தத் திறந்த எடுகோளிலும் $\delta_1$ இன் தனிமாறி
இடம்பெறாமலும் இருந்தால், விளைவாகக் கிடைக்கும்
$\Subst{\delta_1}{\delta_2}{!B^x}$ சரியான !!{proof} ஆகும்.
அதன் திறந்த எடுகோள்கள் $\delta_1$ மற்றும் $\delta_2$ ஆகியவற்றின்
திறந்த எடுகோள்களே. !!{proof}s ஐ ஒட்டும்போது இந்நிபந்தனைகள்
பொதுவாக நிறைவுறும். $\delta_1$, $\delta_2$ ஆகிய இரண்டும்
ஒரே பெரிய !!{proof} இன் உள்நிறுவல்களாக இருந்தால் முதல் நிபந்தனை
நிறைவுறும். இரண்டாவது நிறைவுறாவிட்டால், அது நிறைவுறுமாறு
$\delta_1$ இன் தனிமாறிகளுக்கு வேறு பெயரிடலாம்.

% SEGMENT OLP-0662-S03
\Log{N2c} மற்றும்~\Log{N2i} க்கும் இதற்கொத்த வரையறை பொருந்தும்.

\begin{defn}
$\delta_1$ என்பது $x:!B, \Gamma_1 \Sequent !A$ இன்
\Log{N2c}-!!{proof} (அல்லது \Log{N2i}-!!{proof}) ஆகவும்,
$\delta_2$ என்பது $\Gamma_2 \Sequent !B$ இன்
\Log{N2c}-!!{proof} (அல்லது \Log{N2i}-!!{proof}) ஆகவும் இருந்தால்,
$\Subst{\delta_1}{\delta_2}{x:!B}$ என்பது $\delta_1$ இல்
உள்ள $x:!B \Sequent !B$ வடிவிலான ஒவ்வொரு அடிக்கோளையும்
$\delta_2$ ஆல் பதிலிட்டு, அவ்வடிக்கோள்களுக்குக் கீழுள்ள
ஒவ்வொரு தொடரணியின் சூழலிலும் $\Gamma_2$ ஐச் சேர்ப்பதன் விளைவு.
\end{defn}

% SEGMENT OLP-0662-S04
\begin{prop}
$\delta_1 \Proves x:!B, \Gamma_1 \Sequent
!A$ மற்றும் $\delta_2 \Proves \Gamma_2 \Sequent !B$ என்றால்,
$\delta_1$ இன் எந்தத் தனிமாறியும் $\Gamma_2$ இல்
இடம்பெறாதபோது,
$\Subst{\delta_1}{\delta_2}{x:!B} \Proves \Gamma_1, \Gamma_2 \Sequent
!A$.
\end{prop}

\begin{proof}
  $\pheight{\delta_1}$ மீது தொகுத்தறிதலால் நிறுவலாம்.
\end{proof}

\end{document}
